% Notebook 22 quantum-activity transport module. No preamble. Prefix qt-.
\section{Strict-support transport of analytic quantum activities}
\label{qt-section}

We now substitute the actual stationary current background into a
supplied family of quantum activities. This is a bounded finite-step
map. Its input has only one fixed positive support-cardinality reserve;
we do not assume an arbitrary reserve for quantum activities. The
number of Taylor coordinates is proportional to the original carrier
size, and we explicitly pay that cost. Constants are uniform in the
old refinement $M$ and compatible periodic volume, with fixed $L$.

\subsection{Carrier geometry and the precise analytic inputs}
\label{qt-input-section}

Let $\Lambda$ be the input cell graph and $\Lambda'$ the next field's
cell graph. A projection $\pi:\Lambda\to\Lambda'$ has fibers of size
at most $F$ and sends every nonempty connected input carrier to a
connected set. For ordinary $L$-cube blocking, $F=L^d$. Identity
projection is allowed. Write $d'$ for graph distance on $\Lambda'$ and
define
\begin{equation}\label{qt-distance-sum}
 S_\delta=\sup_{z\in\Lambda'}\sum_{j\in\Lambda'}e^{-\delta d'(z,j)}
 \le\sum_{j\in\mathbb Z^d}e^{-\delta|j|_1}<\infty,
 \qquad \delta>0.
\end{equation}
The bound is uniform on compatible tori: choose one minimal-distance
integer representative of every periodic site. Short periodic
directions do not enlarge the sum.

Each input cell carries at most $D$ scalar complex background
coordinates. For a nonempty connected $H\subset\Lambda$, suppose
$E_H(a|_H)$ is bounded and holomorphic on
\begin{equation}\label{qt-input-domain}
 |a_i^\alpha|<b,\qquad i\in H,\quad 1\le\alpha\le D,
\end{equation}
and put $e_H=\norm{E_H}_b$. The only activity hypothesis is
\begin{equation}\label{qt-input-norm}
 \sup_{i\in\Lambda}\sum_{H\ni i}e^{a_0|H|}e_H\le\epsilon,
 \qquad a_0>0\text{ fixed}.
\end{equation}
Here $D$ counts background scalar coordinates, not automatically
retained Gaussian coordinates. A product of Euclidean fiber balls of
radius $b_E$ contains the scalar polydisc of radius $b_E/\sqrt D$;
make this fixed norm conversion before applying the formulas below.

The actual stationary potential on the next-field product domain of
radius $r$ has the strict support expansion
\begin{align}
 a_i^\alpha(w)&=\sum_{\varnothing\ne X\subset\Lambda'}
                   u_{i,\alpha;X}(w|_X),
 &u_{i,\alpha;X}(0)&=0,                                  \label{qt-background-series}\\
 \sup_{i,\alpha}\sum_X
 e^{\kappa\tau(\{\pi(i)\}\cup X)}
       \norm{u_{i,\alpha;X}}_r&\le s(r)\le C_a r,
 &\kappa&>0.                                             \label{qt-background-norm}
\end{align}
Here $\tau$ is graph Steiner length, allowing auxiliary vertices.
The preceding stationary-background theorem supplies this estimate;
its stronger support-cardinality weight is not needed here. A
maximum-coordinate bound on its full current-field block fiber
implies \eqref{qt-background-norm} for each scalar component, without
summing over the fiber.

If only the total $a_i(0)$ vanishes, subtract each summand's constant.
Absolute convergence preserves the total and increases $s$ by at most
a factor two. Incorporate that factor into $C_a$. For a nonzero
reference potential, center both \eqref{qt-input-domain} and
\eqref{qt-background-series} at that reference and use the available
distance to the analytic boundary in place of $b$.

\begin{remark}[The present physical indexing]\label{qt-indexing}
The normalized Gaussian activities supplied by Notebooks 18 and 21
are already indexed by $L$-blocks of the current unit lattice. The
next-coarse $w$ labels are the same block index graph. The present
physical substitution therefore uses
\[
 \pi=\mathrm{id},\qquad F=1,\qquad
 D=D_{\mathrm{bg}}=dL^d\dim\mathfrak g.
\]
The last dimension includes all current background links assigned to
the block; it differs from the retained Gaussian dimension
$d(L^d-1)\dim\mathfrak g$. After substitution, these labels are unit
cells of the next current lattice. A further projection for the
following integration is a distinct operation, discussed below.
It must not be counted as an additional gain during this substitution.
\end{remark}

\subsection{Bounded transport with an explicit loss budget}

Choose $a_{\mathrm{out}},\mu\ge0$ and $\delta>0$, and set
\begin{equation}\label{qt-budget-notation}
 \beta=a_{\mathrm{out}}+\mu,\qquad
 t=\frac{s(r)}b<1,\qquad
 \alpha_D(t)=D\log\frac1{1-t}.
\end{equation}
For a connected output carrier $K$, let $\ell(K)$ denote metric
minimum spanning-tree length. Define
\begin{equation}\label{qt-output-norm}
 \norm{E'}_{a_{\mathrm{out}},\mu;r}
 =\sup_{z\in\Lambda'}\sum_{K\ni z}
      e^{a_{\mathrm{out}}|K|+\mu\ell(K)}\norm{E'_K}_r.
\end{equation}

\begin{theorem}[Finite-step activity transport]\label{qt-theorem}
Assume \eqref{qt-input-norm}--\eqref{qt-background-norm} and
\begin{equation}\label{qt-budget}
 \beta+\delta\le\kappa,\qquad
 \beta+\alpha_D(t)\le a_0.
\end{equation}
There are strictly supported holomorphic activities $E'_K(w|_K)$
on nonempty connected $K\subset\Lambda'$, with $E'_K(0)=0$, such that
\begin{align}
 \sum_H E_H(a(w)|_H)
   &=\mathrm{Vac}_\Lambda+\sum_K E'_K(w|_K),
 &\mathrm{Vac}_\Lambda&=\sum_H E_H(0),                 \label{qt-exact-identity}\\
 \norm{E'}_{a_{\mathrm{out}},\mu;r}
   &\le F S_\delta\epsilon.                            \label{qt-transport-bound}
\end{align}
All regroupings converge absolutely. The bound is independent of
$M$ and finite periodic volume. It is not a contraction assertion.
\end{theorem}

A concrete choice with strictly positive output weights is
\begin{align}
 \beta&=\min\{a_0/4,\kappa/4\},
 &a_{\mathrm{out}}&=\mu=\beta/2,
 &\delta&=\kappa/4,\notag\\
 \frac{s(r)}b&\le\min\{1/2,a_0/(8D)\}.                 \label{qt-concrete-budget}
\end{align}
Indeed $\alpha_D(t)\le2Dt\le a_0/4$. Thus a sufficient radius is
\begin{equation}\label{qt-radius}
 r\le\frac b{C_a}\min\{1/2,a_0/(8D)\},
\end{equation}
together with the stationary-background theorem's own radius
restriction. All centering and norm-conversion constants are included
in $b$ or $C_a$. Nothing requires $C_a<1$.

\begin{proof}
Expand each activity in its actual $D|H|$ scalar coordinates:
\begin{equation}\label{qt-taylor}
 E_H(a|_H)=\sum_{\nu\in\mathbb N^{D|H|}}c_{H,\nu}a^\nu,
 \qquad |c_{H,\nu}|\le e_H b^{-|\nu|}.
\end{equation}
Coordinatewise Cauchy on smaller polydiscs, followed by a limit,
gives the coefficient bound. There is no dimension-free absolute
Taylor norm being assumed.

For every scalar occurrence in a monomial $a^\nu$, insert
\eqref{qt-background-series}. Its chosen support $X$ inherits the
root $\pi(i)$ of that scalar coordinate. For each such rooted set,
fix a deterministic minimizing connected tree $T(i,X)$ containing
$\{\pi(i)\}\cup X$. Ties can be resolved by any fixed finite-volume
rule. For a decorated term define
\begin{equation}\label{qt-hull}
 K=\pi(H)\cup\bigcup_{\text{occurrences}}\mathrm{Vert}\,T(i,X),
 \qquad T_{\mathrm{tot}}
       =\sum_{\text{occurrences}}\tau(\{\pi(i)\}\cup X).
\end{equation}
The original projected carrier is retained even if some of its
variables are absent from the monomial. It is connected, and every
attached tree has a root in it. Consequently
\begin{equation}\label{qt-hull-geometry}
 |K|\le|\pi(H)|+T_{\mathrm{tot}},\qquad
 z\in K\ \Longrightarrow\ d'(z,\pi(H))\le T_{\mathrm{tot}}.
\end{equation}
A tree with $q$ edges adds at most $q$ vertices after attachment at
its root. A vertex on an attached tree is at distance at most that
tree's length from $\pi(H)$. Overlaps and periodic wrapping can
only improve these inequalities.

The decorated function reads only $w|_K$. Its true variable set may
be smaller: $K$ is a declared connected carrier, not a claim of
minimal dependence. A nonconstant monomial is a product of factors
vanishing at zero and hence also vanishes at zero. Define $E'_K$ as
the sum of these nonconstant terms having carrier $K$. Set aside
the constant Taylor terms as $\mathrm{Vac}_\Lambda$.

For every connected lattice vertex set $K$,
\begin{equation}\label{qt-connected-mst}
 \ell(K)=|K|-1,
 \qquad e^{a_{\mathrm{out}}|K|+\mu\ell(K)}\le e^{\beta|K|}.
\end{equation}
Indeed the induced connected graph has a unit-edge spanning tree,
while every metric tree has $|K|-1$ edges of distance at least one.
Graph Steiner length is used for the attached rooted sets because
deleting intermediate roots from a terminal metric spanning tree
need not decrease its length. At the output all connecting vertices
are retained, so \eqref{qt-connected-mst} applies exactly.

Put $h=|\pi(H)|$ and $k=|H|$. For a fixed output anchor $z$,
\eqref{qt-hull-geometry} gives
\begin{align}
 \mathbf1_{\{z\in K\}}
 e^{a_{\mathrm{out}}|K|+\mu\ell(K)}
 &\le e^{\beta h}e^{-\delta d'(z,\pi(H))}
       \prod_{\text{occurrences}}
       e^{(\beta+\delta)\tau(\{\pi(i)\}\cup X)}.       \label{qt-anchor-hull}
\end{align}
If $z\notin K$ this is automatic; otherwise insert
$1\le e^{\delta(T_{\mathrm{tot}}-d'(z,\pi(H)))}$.
The first condition in \eqref{qt-budget} lets us sum each scalar
occurrence's support using \eqref{qt-background-norm}, at cost $s(r)$.
The exact sum over scalar multiindices is
\begin{equation}\label{qt-coordinate-count}
 \sum_{\nu\in\mathbb N^{Dk}}t^{|\nu|}
    =(1-t)^{-Dk}=e^{\alpha_D(t)k}.
\end{equation}
Omitting $\nu=0$ replaces the last expression by
$e^{\alpha_D(t)k}-1$. Thus the full contribution from a fixed $H$
is bounded by
\begin{equation}\label{qt-fixed-carrier-bound}
 e_H e^{\beta h+\alpha_D(t)k}
                  e^{-\delta d'(z,\pi(H))}.
\end{equation}
Since $h\le k$, the second condition in \eqref{qt-budget} pays
the two exponential costs from the given reserve $a_0 k$.
Finally,
\begin{equation}\label{qt-reroot}
 e^{-\delta d'(z,\pi(H))}
    \le\sum_{i\in H}e^{-\delta d'(z,\pi(i))}.
\end{equation}
Summing \eqref{qt-fixed-carrier-bound} over $H$ therefore gives
\begin{align}
 \sum_i e^{-\delta d'(z,\pi(i))}
             \sum_{H\ni i}e^{a_0|H|}e_H
 &\le F S_\delta\epsilon,                              \label{qt-final-majorant}
\end{align}
which proves \eqref{qt-transport-bound}.

This positive majorant also justifies every substitution and
regrouping. At finite volume, summing over output anchors bounds
the total absolute series. Normal convergence on smaller product
polydiscs proves holomorphy. No order sum has been interchanged
with an output supremum: for each fixed $H$ the scalar series is
summed exactly in \eqref{qt-coordinate-count}, before taking the
anchored carrier sum. In particular, no crude $|H|^n$ slot count
with a missing factorial is hidden in the argument.
\end{proof}

\subsection{Vacuum and the smaller-radius improvement}

The vacuum has the explicit density assignment
\begin{equation}\label{qt-vacuum-density}
 v_i=\sum_{H\ni i}\frac{E_H(0)}{|H|},\qquad
 \mathrm{Vac}_\Lambda=\sum_i v_i,
 \qquad |v_i|\le\epsilon,
 \qquad |\mathrm{Vac}_\Lambda|\le|\Lambda|\epsilon.
\end{equation}
Grouping these densities into a block costs at most $F$.

\begin{corollary}[Vacuum-subtracted radius bound]\label{qt-small-radius}
If the hypotheses of Theorem~\ref{qt-theorem} hold with
\begin{equation}\label{qt-spare-reserve}
 \rho=a_0-\beta-\alpha_D(t)>0,
\end{equation}
then also
\begin{equation}\label{qt-small-radius-bound}
 \norm{E'}_{a_{\mathrm{out}},\mu;r}
 \le F S_\delta\epsilon\,\alpha_D(t) C_\rho,
 \qquad C_\rho=\sup_{k\ge1}k e^{-\rho k}\le\frac1{e\rho}.
\end{equation}
\end{corollary}
\begin{proof}
For nonconstant terms use
$e^{\alpha_D(t)k}-1\le\alpha_D(t)k e^{\alpha_D(t)k}$
in place of \eqref{qt-coordinate-count}. The spare
$e^{-\rho k}$ pays the factor $k$. The remaining rerooting sum
is unchanged.
\end{proof}
One may take the smaller of \eqref{qt-transport-bound} and
\eqref{qt-small-radius-bound}. This is improved smallness on a
strictly smaller field domain. It does not prove a contraction
at a common radius under repeated RG steps.

\subsection{An optional connected-carrier compression estimate}
\label{qt-compression-section}

\begin{lemma}[Connected carriers under cube projection]\label{qt-compression}
For odd $L\ge5$, put
\begin{equation}\label{qt-compression-constants}
 r_L=(L-1)/2,\qquad s_L=\lfloor(r_L-1)/2\rfloor,
 \qquad k_L=s_L+1.
\end{equation}
If $H$ is connected and $\pi_L$ is $L$-cube projection, then
\begin{equation}\label{qt-compression-bound}
 |\pi_L(H)|\le2^d\max\{1,|H|/k_L\}
           \le2^d+(2^d/k_L)|H|.
\end{equation}
The estimate holds on compatible rectangular tori.
\end{lemma}
\begin{proof}
Choose a maximal subset of $H$ separated by at least $r_L$ in
the periodic supremum metric. The supremum balls of radius
$r_L-1$ around its $J$ centers cover $H$. Each such ball has
side length $2r_L-1<L$ and meets at most $2^d$ coarse blocks,
so $|\pi_L(H)|\le2^d J$. If $J=1$ the result follows.
If $J\ge2$, the graph-distance balls of radius $s_L$ about
the centers are disjoint, since $2s_L<r_L$. For each center,
a path in $H$ to another center contains at least $s_L+1$
distinct vertices before leaving its ball. Hence
$Jk_L\le|H|$. Identifications in short periodic directions
only reduce the block counts; the same separation and path
argument uses periodic distances.
\end{proof}

\begin{corollary}[Transport with a proved carrier gain]\label{qt-compressed-transport}
Suppose the projection in Theorem~\ref{qt-theorem} satisfies
$|\pi(H)|\le c_0+c_1|H|$ on connected carriers. The second
condition in \eqref{qt-budget} may be replaced by
\begin{equation}\label{qt-compressed-budget}
 c_1\beta+\alpha_D(t)\le a_0,
\end{equation}
and the resulting bound is
\begin{equation}\label{qt-compressed-bound}
 \norm{E'}_{a_{\mathrm{out}},\mu;r}
       \le e^{\beta c_0}F S_\delta\epsilon.
\end{equation}
The smaller-radius estimate also holds with the factor
$e^{\beta c_0}$ and
$\rho=a_0-c_1\beta-\alpha_D(t)>0$.
\end{corollary}
\begin{proof}
Keep $h=|\pi(H)|$ in \eqref{qt-fixed-carrier-bound}, and
use $\beta h\le\beta c_0+c_1\beta |H|$. All other parts
of the proof are unchanged.
\end{proof}
Lemma~\ref{qt-compression} provides
$c_0=2^d$, $c_1=2^d/k_L$. A sufficiently large fixed $L$
can therefore help pay stronger cardinality weights. But the
attached background trees must still obey
$\beta+\delta\le\kappa$ in the graph used by the substitution.
A larger spatial exponent cannot be assumed without a separate
regrouping proof and its accompanying cardinality and domain costs.
The constant $e^{\beta2^d}L^d S_\delta$ can be large.

In the actual convention of Remark~\ref{qt-indexing}, current
substitution uses identity projection. Afterwards, its connected
carriers are unit-cell carriers of the next current field $w$.
Regrouping them into the following integration's $L$-block Gaussian
cells is a legitimate new projection. It must include that later
pivot-chart enlargement and background-domain budget; it is not
a second projection hidden in \eqref{qt-exact-identity}.

\subsection{What the estimate does and does not establish}

Connectedness of the input carrier is necessary for this
cardinality-only hypothesis. Without it, an activity reading two
sites at arbitrarily large separation can have fixed cardinality
norm, while any connected output carrier containing both has size
at least their distance plus one. A spatial input norm could pay
this cost, but it is not part of \eqref{qt-input-norm}.

There is no automatic contraction: a constant singleton activity
$E_{\{i\}}=c$ at each old cell gives $L^d c$ after summation in
a full coarse block. Even after vacuum subtraction, an identity
background map and a linear singleton activity have no strict
contraction at the same analytic radius. The gain in
\eqref{qt-small-radius-bound} is explicitly a radius loss.

For example, a Gaussian output reserve $a_0\le1/2$ cannot simply
be re-used as an interaction input demanding exponent $4$.
One must verify a compression/enlargement budget, a suitable
stationary-background exponent, and the next domain. The
stationary branch has a one-step radius with an $L$-dependent
constant; no supremum-domain contraction by $L^{-1}$ has been
proved. These facts leave a genuine induction interface.

Every transported coefficient is strictly a function of its
declared $w$ variables. Its fixed baseline may still depend on
$M$ and the finite regulator, and a deterministic choice of
connecting trees need not give coefficientwise volume consistency.
The uniform bounds alone imply no thermodynamic consistency.

Finally, the theorem transports supplied quantum activities
through the matched classical stationary background. It does
not identify the full interacting quantum stationary point with
that classical branch, absorb inherited interactions into a new
minimization, or control running couplings, large-field histories,
or observable insertions. It supplies an exact finite-step bridge
whose coefficient, support, volume, and analytic-radius costs
are all displayed.
